\documentclass[10pt,a4paper]{amsart}
\usepackage[margin=19mm]{geometry}
\usepackage{amsmath,amssymb,mathtools}
\usepackage[scr=rsfs,cal=euler]{mathalfa}
\usepackage{xfrac}
\usepackage[unicode,hidelinks,bookmarks=false]{hyperref}
\newcommand{\ie}{i.e.\ }
\newcommand{\cf}{cf.\ }
\newcommand{\resp}{respectively\ }
\newcommand{\Subsection}{\subsection}
\providecommand{\nobreakdash}{\nobreak\hbox{-}}
\setlength{\emergencystretch}{2em}
\multlinegap0pt
\usepackage{xcolor}
\usepackage{amssymb}
\usepackage[shortlabels]{enumitem}
\newtheorem{lem}{Lemma}
\providecommand{\Aut}{{\operatorname{Aut}}}
\providecommand{\C}{{\mathrm{C}}}
\providecommand{\gen}[1]{\langle #1 \rangle} % CHECK: duplicate definition; repeated again below.
\providecommand{\olG}{{\overline{G}}}
\providecommand{\olH}{{\overline{H}}}
\providecommand{\olb}{{\overline{b}}}
\providecommand{\olc}{{\overline{c}}}
\title{On $\{2,3,5\}$-groups with conjugacy classes of distinct sizes}
\author{Wei Zhou, Ilya Gorshkov}
\date{Source: arXiv:2606.22244v1 (CC BY 4.0)}
\begin{document}
\maketitle

\noindent\emph{Source and licence.} On $\{2,3,5\}$-groups with conjugacy classes of distinct sizes. Version arXiv:2606.22244v1 (CC BY 4.0), distributed by arXiv under the Creative Commons Attribution 4.0 International licence, http://creativecommons.org/licenses/by/4.0/, as recorded in the arXiv licence metadata for this exact version and shown on the abstract page https://arxiv.org/abs/2606.22244v1. Full licence notice: \texttt{licenses/arxiv-2606.22244-CC-BY-4.0.txt}.\par\smallskip

\noindent\textit{Editorial scope.} Counted unit: the article's lem real3auta6 (source0.tex lines 1166--1170). The article's complete original proof of it is delivered with it (source0.tex lines 1172--1211, one \texttt{proof} environment of 40 lines). No mathematics is written, repaired or re-derived: the statement and the proof are the article's own lines, in the article's own notation, and no retained line is substituted or re-flowed.  No further source-local context is needed: every cross-reference of every retained line resolves inside the two delivered spans.  Nothing was omitted from any retained span, so the package carries no omission record.  No retained line contains a citation, so the wrapper carries no bibliography and no bibliography entry.  No retained line contains a figure environment, an image-inclusion command or drawing code, so no image is delivered and the package declares no asset dependency.  Editorial formatting only, in addition: an AMS \texttt{amsart} preamble that restates the article's own declarations and macros used by the retained lines, namely the environments \texttt{lem} on separate counters, together with the standard packages those lines need.  The retained environments print in this document as Lemma 1 in order of appearance, whereas the article prints the same items with the numbering of its own full document.  The complete native source of the article as distributed in the arXiv e-print for this exact version is bundled unchanged as \texttt{sources/203335/source0.tex}.
\begin{lem}
\label{real3auta6}
    Let $G$ be a finite group and $N$ be a normal subgroup of $G$.
    If $|N|=3$ and $G/N\cong S_6$ or $\Aut(A_6)$, then $G$ is not rational.
\end{lem}

\begin{proof}
    Suppose that the assertion of this lemma is false.
    Let $N=\gen{a}$, where $a$ is an element of order $3$,
    and set $\olG=G/N$.
    Since $G$ is rational, we have $G/\C_G(a)\cong \Aut(C_3)\cong C_2$.
    Let $H =\C_G(a)$. Note that $|\olG:\olH|=2$.

    Assume that $\olG \cong S_6$.
    Hence $\olH \cong A_6$.
    Let $b\in H$ such that $|b|=5$.
    Since the centralizer of a $5$-cycle in $S_6$ is generated by the cycle itself,
    we have $\C_{\olG}(\olb)=\gen{\olb}\leq \olH$.
    Since $(|N|,|b|)=1$, we have $\C_{\olG}(\olb)=\C_G(b)N/N$.
    Hence $\C_G(b)\leq H$.
    Since $G$ is rational, 
    there exists an element $x\in G$ such that $(ab)^x=(ab)^{11}$.
    It follows that $a^x=a^2$ and $b^x=b$.
    The latter implies $x\in \C_G(b)\leq H$, so $a^x=a$, a contradiction.

    Therefore, we must have $\olG \cong \Aut(A_6)$.
    If $\olH \cong S_6$, 
    choose $b\in H$ such that $|b|=2$ and $\olb$ is an odd permutation in $\olH$.
    Thus $\C_{\olG}(\olb)\leq \olH$.
    Similarly, we have $\C_G(b)\leq H$.
    Since $G$ is rational, 
    there exists an element $x$ in $G$ such that $(ab)^x=(ab)^5$.
    We have $a^x=a^2$ and $b^x=b$, a contradiction.
    
    Hence $\olH$ is not isomorphic to $S_6$.
    Hence $\olH$ is isomorphic to either $M_{10}$ or $PGL_2(9)$.
    In these cases, there exists an element $c$ in $H$ such that $|c|=8$.
    Since every element of order $8$ in $\Aut(A_6)$ is self-centralizing,
    we have $\C_{\olG}(\olc)=\gen{\olc}\leq \olH$.
    Similarly, we have $\C_G(c)\leq H$.
    Since $G$ is rational, 
    there exists an element $y\in G$ such that $(ac)^y = (ac)^{17}$.
    It follows that $a^y=a^2$ and $c^y=c$, a contradiction.
    
    Hence $G$ is not rational.
\end{proof}

\end{document}
